\documentclass[11pt,a4paper]{article}

\usepackage[utf8]{inputenc}
\usepackage[T1]{fontenc}
\usepackage[ngerman]{babel}
\usepackage{lmodern}
\usepackage{geometry}
\usepackage{xcolor}
\usepackage{graphicx}
\usepackage{tabularx}
\usepackage{longtable}
\usepackage{booktabs}
\usepackage{array}
\usepackage{enumitem}
\usepackage{hyperref}
\usepackage{listings}
\usepackage{tikz}
\usetikzlibrary{arrows.meta,positioning,shapes.geometric}

\geometry{margin=24mm}
\hypersetup{colorlinks=true,linkcolor=blue!50!black,urlcolor=blue!50!black}
\setlist[itemize]{leftmargin=*,itemsep=2pt}
\setlist[enumerate]{leftmargin=*,itemsep=2pt}
\renewcommand{\arraystretch}{1.18}

\definecolor{brand}{HTML}{00684F}
\definecolor{accent}{HTML}{C11574}
\definecolor{soft}{HTML}{F5F7FB}
\definecolor{line}{HTML}{D0D5DD}
\definecolor{textmuted}{HTML}{475467}

\lstdefinelanguage{PlantUML}{
  morekeywords={@startuml,@enduml,actor,participant,database,cloud,node,artifact,package,rectangle,usecase,entity,skinparam,hide,left,to,right,direction,start,stop,if,then,else,endif,repeat,while},
  sensitive=false,
  morecomment=[l]{'},
}
\lstset{
  basicstyle=\ttfamily\small,
  keywordstyle=\color{brand}\bfseries,
  commentstyle=\color{textmuted},
  breaklines=true,
  frame=single,
  rulecolor=\color{line},
  columns=fullflexible
}

\newcommand{\docTitle}{Anforderungsspezifikation StudentGo}
\newcommand{\version}{1.0}
\newcommand{\status}{Entwurf}
\newcommand{\product}{StudentGo}

\newcolumntype{Y}{>{\raggedright\arraybackslash}X}
\newcommand{\priority}[1]{\textbf{#1}}
\newcommand{\must}{\priority{Muss}}
\newcommand{\should}{\priority{Soll}}
\newcommand{\could}{\priority{Kann}}

\begin{document}

\begin{titlepage}
  \centering
  \vspace*{18mm}
  {\Huge\bfseries \docTitle\par}
  \vspace{8mm}
  {\Large Mobile Campus-App fuer Studierende der HFU\par}
  \vfill
  \begin{tabular}{ll}
    Produkt & \product\\
    Version & \version\\
    Status & \status\\
    Datum & \today\\
    Bezug & Expo/React Native App mit Express/Prisma Backend
  \end{tabular}
  \vfill
  {\color{brand}\rule{0.7\textwidth}{1.2pt}}
\end{titlepage}

\tableofcontents
\newpage

\section{Zweck und Geltungsbereich}
Diese Anforderungsspezifikation beschreibt die fachlichen und technischen Anforderungen an \product. Die App buendelt studienrelevante Alltagsfunktionen: Stundenplan, Mensa, Kontakte, To-dos, private Notizen, Rollenverwaltung und kontrollierte Inhaltsaenderungen. Sie basiert auf der bestehenden Codebasis mit Expo Router, React Native, Express, Prisma und SQLite.

\subsection{Ziele}
\begin{itemize}
  \item Studierende erhalten einen kompakten, mobilen Zugriff auf persoenliche und oeffentliche Campus-Informationen.
  \item Managerinnen und Manager koennen Aenderungen fuer verwaltete Inhalte vorbereiten, ohne direkt produktive Daten zu ueberschreiben.
  \item Administratorinnen und Administratoren pruefen Antraege, verwalten Benutzer und importieren externe Datenquellen.
  \item Kerninformationen bleiben im Offline-Fall lesbar; neue persoenliche Eintraege werden nach Verbindungswiederkehr synchronisiert.
\end{itemize}

\subsection{Systemkontext}
\product{} besteht aus einer mobilen bzw. webbasierten Expo-App, einem Express-Backend und einer SQLite-Datenbank. Externe Systeme sind der HFU-StarPlan fuer Stundenplandaten und der SWFR-Speiseplan fuer Mensadaten.

\section{Stakeholder und Rollen}
\begin{tabularx}{\textwidth}{@{}lY@{}}
\toprule
Rolle & Interesse\\
\midrule
Student:in & Schneller Zugriff auf Stundenplan, Aufgaben, Notizen, Kontakte, Mensa und geteilte Veranstaltungen.\\
Manager:in & Erstellung von Aenderungsantraegen fuer Kontakte, Mensa, Lehrveranstaltungen und Studieninformationen.\\
Administrator:in & Benutzerverwaltung, Rollensteuerung, Freigabe oder Ablehnung von Aenderungen, Importsteuerung.\\
Backend-Betreiber:in & Sicherer und wartbarer Betrieb der API, Datenbank und Importprozesse.\\
\bottomrule
\end{tabularx}

\section{Funktions-Master}
\begin{longtable}{@{}p{18mm}p{35mm}p{77mm}p{18mm}@{}}
\toprule
ID & Funktion & Beschreibung & Prioritaet\\
\midrule
\endhead
FM-01 & Registrierung und Login & Benutzer registrieren sich mit Name, E-Mail, Passwort und Rolle; Accounts werden ueber Bestaetigungstoken aktiviert. & \must\\
FM-02 & Rollenbasierter Zugriff & Sessions enthalten eine aktive Rolle; Admin darf alle Rollen nutzen, Manager keine Admin-Rolle, User nur User. & \must\\
FM-03 & Benutzerverwaltung & Admins legen Benutzer an, aendern Rollen, aktualisieren Accounts und loeschen Benutzer. & \must\\
FM-04 & Aenderungsworkflow & Manager und Admins erstellen Aenderungsantraege; Admins genehmigen oder lehnen diese ab. & \must\\
FM-05 & Kontakte & Oeffentliche und persoenliche Kontakte anzeigen, suchen und persoenliche Kontakte anlegen. & \must\\
FM-06 & To-dos & Persoenliche To-dos mit Unteraufgaben anlegen, anzeigen und abschliessen. & \must\\
FM-07 & Private Notizen & Studieninformationen und persoenliche Notizen anzeigen; private Inhalte koennen clientseitig verschluesselt werden. & \should\\
FM-08 & Stundenplan & Wochenplan anzeigen, manuelle Veranstaltungen pflegen, Besuche markieren und Modul-Praeferenzen speichern. & \must\\
FM-09 & StarPlan-Import & Importoptionen laden, Kurse auswaehlen und Veranstaltungen aus HFU-StarPlan uebernehmen. & \should\\
FM-10 & Einladungen & Veranstaltungen mit anderen Benutzern teilen; Empfaenger koennen annehmen oder ablehnen. & \should\\
FM-11 & Mensa & Mensen und Speiseplaene anzeigen; Admin kann SWFR-Import ausloesen. & \should\\
FM-12 & Offline-Faehigkeit & GET-Antworten cachen, ausgewaehlte POST-Aktionen lokal vormerken und spaeter wiederholen. & \must\\
FM-13 & Theme und Bedienung & Helles/dunkles Erscheinungsbild, Statusanzeigen und Swipe-Navigation bereitstellen. & \could\\
\bottomrule
\end{longtable}

\section{User-Stories}
\begin{longtable}{@{}p{18mm}p{105mm}p{26mm}@{}}
\toprule
ID & Story & Akzeptanzkriterien\\
\midrule
\endhead
US-01 & Als Student:in moechte ich mich registrieren und bestaetigen, damit ich persoenliche Daten speichern kann. & Bestaetigungstoken wird erzeugt; vor Bestaetigung existiert nur ein PendingAccount; nach Bestaetigung ist Login moeglich.\\
US-02 & Als Student:in moechte ich meine Kontakte durchsuchen, damit ich Personen schnell finde. & Suche filtert Name, Rolle, E-Mail und Raum; persoenliche Kontakte werden markiert.\\
US-03 & Als Student:in moechte ich To-dos mit Unteraufgaben erfassen, damit ich Studienaufgaben strukturieren kann. & To-do wird dem Account zugeordnet; Unteraufgaben sind separat abschliessbar; erledigte Aufgaben werden erkennbar getrennt.\\
US-04 & Als Student:in moechte ich private Notizen verschluesseln, damit sensible Inhalte nur mit meinem Private Key lesbar sind. & Public Key wird am Account gespeichert; verschluesselte Payloads enthalten Key und IV; ohne passenden Private Key wird ein Hinweis angezeigt.\\
US-05 & Als Student:in moechte ich meinen Stundenplan importieren, damit ich ihn nicht manuell uebertragen muss. & Importoptionen werden geladen; Auswahl erzeugt Lessons und ScheduleDays; Duplikate werden anhand stabiler Modulschluessel vermieden.\\
US-06 & Als Student:in moechte ich eine Veranstaltung teilen, damit Kommiliton:innen sie uebernehmen koennen. & Einladung referenziert Lesson, Sender und Empfaenger; Status ist initial PENDING; Annahme/Ablehnung wird gespeichert.\\
US-07 & Als Manager:in moechte ich Aenderungen beantragen, damit fachliche Inhalte kontrolliert gepflegt werden. & Antrag enthaelt Entity, Action, Payload und optional TargetId; Status ist PENDING; eigene Antraege sind sichtbar.\\
US-08 & Als Admin moechte ich Aenderungsantraege freigeben, damit produktive Daten erst nach Pruefung geaendert werden. & Approve wendet Payload auf Zielentitaet an; Reject speichert Review-Notiz; Status und Review-Zeitpunkt werden gesetzt.\\
US-09 & Als Admin moechte ich Benutzerrollen verwalten, damit Verantwortlichkeiten korrekt abgebildet sind. & Rollen USER, MANAGER und ADMIN sind waehlbar; unzulaessige Rollenwechsel werden serverseitig verhindert.\\
US-10 & Als Nutzer:in moechte ich bei Netzproblemen weiterarbeiten, damit kurze Verbindungsabbrueche mich nicht blockieren. & Cache liefert bekannte GET-Daten; lokale Creates erhalten Pending-Status; SyncContext spielt Requests spaeter ab.\\
\bottomrule
\end{longtable}

\section{Nichtfunktionale Anforderungen}
\begin{longtable}{@{}p{20mm}p{35mm}p{86mm}@{}}
\toprule
ID & Kategorie & Anforderung\\
\midrule
\endhead
NF-01 & Sicherheit & Passwoerter werden mit Salt und scrypt gehasht; Sessiontokens werden nur als SHA-256 Hash gespeichert.\\
NF-02 & Datenschutz & Private Schluessel verbleiben clientseitig; verschluesselte Inhalte werden als Payload, Key und IV gespeichert.\\
NF-03 & Zugriffsschutz & Admin- und Manager-Endpunkte muessen serverseitig ueber Middleware geschuetzt sein.\\
NF-04 & Offline & Cache und Queue duerfen Benutzerkontext nicht vermischen; Cache-Key beruecksichtigt URL und Authorization-Header.\\
NF-05 & Wartbarkeit & Fachliche Modelle werden zentral in Prisma beschrieben; API und App verwenden klar benannte Module und Contexts.\\
NF-06 & Portabilitaet & App laeuft auf Expo-Zielplattformen Android, iOS und Web; Backend laeuft lokal oder auf Node.js Hosts.\\
NF-07 & Nachvollziehbarkeit & Aenderungsantraege speichern Ersteller, Review-Daten, Status und Payload.\\
NF-08 & Importrobustheit & Externe Importe muessen fehlschlagen duerfen, ohne bestehende App-Funktionen zu blockieren.\\
\bottomrule
\end{longtable}

\section{4+1 Sichtenmodell}
\subsection{Use-Case-Sicht}
Die Use-Case-Sicht beschreibt die wichtigsten Interaktionen zwischen Rollen und System.

\begin{center}
\begin{tikzpicture}[node distance=12mm, actor/.style={draw,rounded corners,fill=soft,minimum width=24mm,minimum height=8mm}, uc/.style={draw,ellipse,fill=white,minimum width=35mm,minimum height=9mm}, line/.style={-Latex}]
  \node[actor] (student) {Student:in};
  \node[actor,below=18mm of student] (manager) {Manager:in};
  \node[actor,below=18mm of manager] (admin) {Admin};
  \node[uc,right=28mm of student] (auth) {Auth};
  \node[uc,below=7mm of auth] (contacts) {Kontakte};
  \node[uc,below=7mm of contacts] (schedule) {Stundenplan};
  \node[uc,below=7mm of schedule] (change) {Aenderungsantrag};
  \node[uc,below=7mm of change] (review) {Review};
  \node[uc,right=24mm of schedule] (offline) {Offline Sync};
  \draw[line] (student) -- (auth);
  \draw[line] (student) -- (contacts);
  \draw[line] (student) -- (schedule);
  \draw[line] (student) -- (offline);
  \draw[line] (manager) -- (change);
  \draw[line] (admin) -- (review);
  \draw[line] (admin) -- (auth);
  \draw[line,dashed] (change) -- (review);
\end{tikzpicture}
\end{center}

PlantUML-Quelle: \texttt{diagrams/usecase.puml}.

\subsection{Logische Sicht}
Die logische Sicht trennt App-Kontexte, Backend-Services und Persistenz.

\begin{center}
\begin{tikzpicture}[node distance=10mm, box/.style={draw,rounded corners,fill=soft,minimum width=36mm,minimum height=9mm,align=center}, ext/.style={draw,rounded corners,fill=white,minimum width=30mm,minimum height=9mm,align=center}, line/.style={-Latex}]
  \node[box] (tabs) {Expo Router\\Screens};
  \node[box,below=of tabs] (contexts) {Auth, Offline,\\Sync, Theme};
  \node[box,right=28mm of tabs] (api) {Express\\REST API};
  \node[box,below=of api] (services) {Auth, Review,\\Import Services};
  \node[box,right=28mm of api] (prisma) {Prisma Client};
  \node[box,below=of prisma] (db) {SQLite};
  \node[ext,above=of api] (starplan) {HFU StarPlan};
  \node[ext,above=of prisma] (swfr) {SWFR};
  \draw[line] (tabs) -- (contexts);
  \draw[line] (tabs) -- node[above]{REST/JSON} (api);
  \draw[line] (api) -- (services);
  \draw[line] (api) -- (prisma);
  \draw[line] (prisma) -- (db);
  \draw[line] (services) -- (starplan);
  \draw[line] (services) -- (swfr);
\end{tikzpicture}
\end{center}

PlantUML-Quelle: \texttt{diagrams/logical-components.puml}.

\subsection{Prozesssicht}
Wichtige Laufzeitprozesse sind Login, Offline-Synchronisation, Aenderungsreview, StarPlan-Import und SWFR-Import. Die Offline-Synchronisation puffert schreibende Aktionen fuer Kontakte und To-dos, markiert sie lokal als pending und spielt sie mit Replay-Header erneut ab.

\begin{center}
\begin{tikzpicture}[node distance=9mm, procstep/.style={draw,rounded corners,fill=soft,minimum width=42mm,minimum height=8mm,align=center}, decision/.style={draw,diamond,aspect=2,fill=white,align=center}, line/.style={-Latex}]
  \node[procstep] (create) {Aktion in App};
  \node[decision,below=of create] (online) {Backend\\erreichbar?};
  \node[procstep,left=18mm of online] (save) {Direkt speichern};
  \node[procstep,right=18mm of online] (queue) {Lokal queuen};
  \node[procstep,below=18mm of queue] (replay) {Bei Verbindung\\Replay senden};
  \node[procstep,below=18mm of online] (reload) {Listen neu laden};
  \draw[line] (create) -- (online);
  \draw[line] (online) -- node[above]{ja} (save);
  \draw[line] (online) -- node[above]{nein} (queue);
  \draw[line] (queue) -- (replay);
  \draw[line] (save) |- (reload);
  \draw[line] (replay) |- (reload);
\end{tikzpicture}
\end{center}

PlantUML-Quellen: \texttt{diagrams/process-sync.puml}, \texttt{diagrams/admin-review-sequence.puml}, \texttt{diagrams/scenario-schedule-import.puml}.

\subsection{Entwicklungssicht}
Die Entwicklungssicht folgt der vorhandenen Projektstruktur. Screens liegen in \texttt{app/(tabs)}, gemeinsam genutzte Zustandslogik in \texttt{contexts}, plattformspezifische Persistenz und Kryptografie in \texttt{lib}, Backend-Code in \texttt{backend/src} und Datenmodell/Migrationen in \texttt{prisma}.

\begin{center}
\begin{tikzpicture}[node distance=8mm, pkg/.style={draw,rounded corners,fill=soft,minimum width=35mm,minimum height=9mm,align=center}, line/.style={-Latex}]
  \node[pkg] (app) {app/(tabs)};
  \node[pkg,right=of app] (contexts) {contexts};
  \node[pkg,right=of contexts] (lib) {lib};
  \node[pkg,below=of contexts] (backend) {backend/src};
  \node[pkg,right=of backend] (prisma) {prisma};
  \node[pkg,left=of backend] (components) {components};
  \draw[line] (app) -- (contexts);
  \draw[line] (app) -- (components);
  \draw[line] (contexts) -- (lib);
  \draw[line] (app) -- (backend);
  \draw[line] (backend) -- (prisma);
\end{tikzpicture}
\end{center}

PlantUML-Quelle: \texttt{diagrams/development-view.puml}.

\subsection{Physische Sicht}
Die physische Sicht besteht aus Clientgeraet, Backend-Host, Datenbank und externen Importquellen.

\begin{center}
\begin{tikzpicture}[node distance=12mm, nodebox/.style={draw,rounded corners,fill=soft,minimum width=42mm,minimum height=13mm,align=center}, cloud/.style={draw,ellipse,fill=white,minimum width=34mm,minimum height=10mm,align=center}, line/.style={-Latex}]
  \node[nodebox] (client) {Mobilgeraet/Web\\Expo App + Local Storage};
  \node[nodebox,right=28mm of client] (backend) {Backend Host\\Node.js Express};
  \node[nodebox,right=28mm of backend] (db) {SQLite\\Datenbank};
  \node[cloud,above=of backend] (starplan) {HFU StarPlan};
  \node[cloud,below=of backend] (swfr) {SWFR};
  \draw[line] (client) -- node[above]{HTTPS REST} (backend);
  \draw[line] (backend) -- (db);
  \draw[line] (backend) -- (starplan);
  \draw[line] (backend) -- (swfr);
\end{tikzpicture}
\end{center}

PlantUML-Quelle: \texttt{diagrams/deployment-view.puml}.

\section{Datenanforderungen}
Das zentrale Datenmodell wird durch Prisma beschrieben. Besonders relevant sind \texttt{User}, \texttt{Session}, \texttt{PendingAccount}, \texttt{ManagementChangeRequest}, \texttt{Contact}, \texttt{Todo}, \texttt{StudyInfo}, \texttt{Lesson}, \texttt{LessonInvitation}, \texttt{Canteen} und \texttt{MealPlan}.

\begin{tabularx}{\textwidth}{@{}lY@{}}
\toprule
Entitaet & Zweck\\
\midrule
User & Account, Rolle, Public Keys und Beziehungen zu persoenlichen Daten.\\
Session & Tokenbasierte Anmeldung mit aktiver Rolle und Ablaufzeit.\\
ManagementChangeRequest & Auditierbarer Aenderungsantrag mit Action, Entity, Payload und Review-Daten.\\
Lesson & Einzelne oder wiederkehrende Veranstaltung inklusive Importquelle, Raum, Dozent:in und Verschluesselungsfeldern.\\
ScheduleImportCache & Zwischenspeicher fuer importierte StarPlan-Daten pro Woche und Auswahl.\\
StudyInfo & Oeffentliche Studieninformationen oder persoenliche, optional verschluesselte Notizen.\\
\bottomrule
\end{tabularx}

PlantUML-Quelle: \texttt{diagrams/data-model.puml}.

\section{Schnittstellen}
\begin{longtable}{@{}p{45mm}p{82mm}p{22mm}@{}}
\toprule
Endpoint & Zweck & Schutz\\
\midrule
\endhead
\texttt{POST /api/auth/register} & Registrierung und PendingAccount-Erzeugung. & oeffentlich\\
\texttt{POST/GET /api/auth/confirm} & Accountbestaetigung per Token. & oeffentlich\\
\texttt{POST /api/auth/login} & Session erzeugen. & oeffentlich\\
\texttt{GET /api/auth/me} & Aktuellen Benutzer laden. & Token\\
\texttt{GET/POST /api/contacts} & Kontakte lesen und persoenliche Kontakte anlegen. & optional/Token\\
\texttt{GET/POST /api/todos} & Persoenliche To-dos lesen und anlegen. & Token\\
\texttt{GET/POST /api/study-info} & Studieninfos und Notizen lesen/anlegen. & optional/Token\\
\texttt{GET /api/schedule} & Stundenplan fuer Woche/Monat laden. & optional/Token\\
\texttt{POST /api/schedule/import} & StarPlan-Veranstaltungen importieren. & Token\\
\texttt{POST /api/schedule/lessons/:id/visit} & Besuch einer Veranstaltung markieren. & Token\\
\texttt{GET/POST /api/manager/change-requests} & Aenderungsantraege lesen/anlegen. & Manager/Admin\\
\texttt{GET /api/admin/users} & Benutzerliste lesen. & Admin\\
\texttt{POST /api/admin/change-requests/:id/approve} & Antrag genehmigen. & Admin\\
\texttt{POST /api/admin/import/swfr-meals} & Mensa-Import ausloesen. & Admin\\
\bottomrule
\end{longtable}

\section{Abnahmekriterien}
\begin{itemize}
  \item Die App kann mit lokalem Backend starten und \texttt{/health} erfolgreich abfragen.
  \item Registrierung, Bestaetigung, Login, Logout und Rollenwahl funktionieren gemaess Zugriffsmatrix.
  \item Admin-Endpunkte sind ohne Admin-Session nicht erreichbar.
  \item Kontakte und To-dos koennen offline vorgemerkt und nach Wiederverbindung synchronisiert werden.
  \item StarPlan- und SWFR-Importfehler fuehren zu sichtbarer Fehlermeldung, aber nicht zum Absturz der App.
  \item Aenderungsantraege koennen erstellt, genehmigt, abgelehnt und im Status nachvollzogen werden.
  \item PlantUML-Quellen liegen unter \texttt{diagrams/} und koennen mit PlantUML in Diagrammartefakte uebersetzt werden.
\end{itemize}

\section{PlantUML-Integration}
Die UML-Quellen liegen als echte PlantUML-Dateien im Ordner \texttt{diagrams/}. Damit LaTeX PlantUML direkt parsen kann, ist zusaetzlich \texttt{plantuml-latex-demo.tex} enthalten. Dieser Build nutzt das lokal installierte LaTeX-Paket \texttt{plantuml}; erforderlich sind \texttt{lualatex -shell-escape} und die Umgebungsvariable \texttt{PLANTUML\_JAR}.

\begin{lstlisting}[language=bash]
export PLANTUML_JAR=/pfad/zu/plantuml.jar
make plantuml-latex
\end{lstlisting}

Der Standard-Build \texttt{make pdf} erzeugt diese Spezifikation ohne externes Jar. Sobald PlantUML installiert ist, erzeugt \texttt{make plantuml} zusaetzlich gerenderte Diagramme aus den \texttt{.puml}-Dateien.

\appendix
\section{PlantUML-Quellen}
\subsection{Use-Case-Diagramm}
\lstinputlisting[language=PlantUML]{diagrams/usecase.puml}
\subsection{Logische Komponenten}
\lstinputlisting[language=PlantUML]{diagrams/logical-components.puml}
\subsection{Offline-Prozess}
\lstinputlisting[language=PlantUML]{diagrams/process-sync.puml}

\end{document}
